% META: {"id":"1SPE-DERGLOBAL-ME-005","chapitre":"1SPE-DERIVATION-GLOBAL","type_objet":"methode","capacites_codes":["C5"],"status":"generated"}
\begin{fichemethode}{M5}{Résoudre un problème d'optimisation}
\quandUtiliser{J'utilise cette méthode lorsque l'énoncé demande de maximiser ou minimiser une grandeur (aire, volume, coût, bénéfice\ldots) soumise à une contrainte.}
\pasApas{\begin{enumerate}
  \item Je modélise : j'exprime la grandeur à optimiser sous la forme d'une fonction $f(x)$ d'une seule variable.
  \item Je détermine le domaine d'étude $I$ en tenant compte des contraintes du problème.
  \item Je vérifie la dérivabilité, puis je calcule $f'(x)$ et je recherche les
        points critiques intérieurs.
  \item J'étudie les variations et je compare les valeurs aux candidats et aux
        bornes incluses pour trouver un extremum global.
  \item Je conclus en répondant à la question posée avec les unités adaptées.
\end{enumerate}}
\exempleRedige{Dans un triangle $ABC$ rectangle en $A$, on a $AB=6$ cm et
$AC=8$ cm, donc $BC=10$ cm. On inscrit un rectangle dont un sommet est $A$,
deux côtés sont portés par $[AB]$ et $[AC]$, et le sommet opposé à $A$ appartient
à $[BC]$. On note $x$ sa largeur le long de $[AB]$ et $h$ sa hauteur, en cm.
\commentaireMarge{Modéliser avec Thalès.}
Les triangles semblables donnent $h/8=(6-x)/6$, donc $h=8-\frac43x$ et
\[f(x)=xh=8x-\frac43x^2.\]
Le domaine des rectangles non dégénérés est $x\in]0\,;\,6[$.
\commentaireMarge{Dériver et étudier le signe.}
On a $f'(x)=8-\frac83x$, positif pour $0<x<3$ et négatif pour $3<x<6$.
L'aire atteint donc son maximum global pour $x=3$ cm ; la hauteur vaut alors
$h=4$ cm et l'aire maximale est $12$ cm².}
\pieges{\begin{itemize}
  \item Oublier de vérifier que le point critique appartient bien au domaine $I$.
  \item Ne pas justifier l'extremum global par les variations et les bornes du domaine.
  \item Confondre la valeur optimale de $x$ avec la valeur optimale de la grandeur $f(x)$.
\end{itemize}}
\verifier{Je m'assure que la réponse est cohérente avec le contexte : une aire ou une longueur ne peut pas être négative.}
\sEntrainer{\refExos{M5}}
\end{fichemethode}

% BEGIN-VERIFY
% from sympy import Matrix, Rational, symbols, expand, diff, solve
% A, B, C = Matrix([0,0]), Matrix([6,0]), Matrix([0,8])
% assert (B-C).dot(B-C) == 100
% x = symbols('x', real=True)
% height = 8-Rational(4, 3)*x
% P = Matrix([x,height])
% # Le sommet opposé du rectangle appartient à la droite BC.
% assert expand((P-B)[0]*(C-B)[1]-(P-B)[1]*(C-B)[0]) == 0
% area = x*height
% assert expand(area-(8*x-Rational(4,3)*x**2)) == 0
% assert solve(diff(area,x),x) == [3]
% assert 0 < 3 < 6
% assert height.subs(x,3) == 4
% assert area.subs(x,3) == 12
% assert expand(12-area-Rational(4,3)*(x-3)**2) == 0
% assert (Rational(4,3)*(x-3)**2).is_nonnegative is True
% END-VERIFY
